% Post-checkpoint-14 theorem. Full author/root/B/C proof audits passed.
\begin{theorem}[Proved: a generic surface detects every original cover]
\label{D:generic-surface-finite-punctures}
Let $k$ be an algebraically closed field of arbitrary characteristic,
let $S$ be a normal proper integral $k$-scheme of dimension $d\ge2$,
and let $T\subset S$ be a finite closed subset. Put $U=S\setminus T$.
There exist an algebraically closed extension $K/k$ of finite
transcendence degree, a normal proper integral surface $X/K$, and a
finite morphism $g:X\to S_K$ onto a closed integral surface, such that
$T_X=g^{-1}(T_K)$ is finite and the composite morphism
\[
       X\setminus T_X\longrightarrow U_K\longrightarrow U
\]
induces a continuous surjection
\[
       \pi_1(X\setminus T_X)\twoheadrightarrow\pi_1(U).
       \tag{D.GS.1}
\]
The extension and surface are fixed independently of the finite cover
of $U$. The target in (D.GS.1) is the fundamental group of the original
$k$-scheme $U$; no surjectivity onto $\pi_1(U_K)$ is asserted.
Consequently, full finite generation for normal proper surfaces with
finite punctures over algebraically closed fields would imply full
finite generation for normal proper schemes of every dimension with
finite punctures.
\end{theorem}

\begin{proof}
For $d=2$ we take $K=k$, $X=S$ and the identity. Suppose $d\ge3$.
Chow's lemma gives a proper modification by a quasi-projective scheme
which is an isomorphism over a dense open. Take the reduced closure of
that open in the modification. Since $S$ is proper, we obtain an
integral projective scheme $Y$ and a proper birational morphism
$q:Y\to S$. Fix a projective embedding $Y\hookrightarrow\mathbf P^N_k$.
There is a nonempty smooth open $W\subset U$ over which $q$ is an
isomorphism; we also regard $W$ as a smooth locally closed subscheme of
$\mathbf P^N_k$. We use only this dense open for the Bertini argument.
In particular, we do not require a resolution of $S$ or normality of a
whole hyperplane section of $Y$
\citep[Tags 0200 and 0FD6]{StacksDGenericSurface}.

We first verify the exact generic-section input, including its application
to a covering. Let $F$ be an algebraically closed field, let
$W_0\hookrightarrow\mathbf P^N_F$ be smooth and integral of dimension
$e\ge3$, and let $V_0\to W_0$ be a connected finite etale cover.
Then $V_0$ is smooth and integral, hence geometrically irreducible.
Pull back the hyperplane linear system to $V_0$. It is basepoint-free.
Choose a closed point $v\in V_0$, with image $w$. Since $V_0\to W_0$
is etale and $W_0\hookrightarrow\mathbf P^N_F$ is an immersion, we can
choose two hyperplanes through $w$ whose local equations have linearly
independent differentials at $v$. Their pulled-back intersection is
smooth of codimension two near $v$. Its irreducible component through
$v$ thus has codimension two in $V_0$. A third hyperplane not passing
through $w$ does not contain that component. These observations verify
conditions (3) and (4) of Stacks, Lemma 37.32.3, as well as its finite-type
and geometric-irreducibility hypotheses. That lemma therefore says that
the pullback of a geometric generic hyperplane is geometrically
irreducible. Meanwhile Lemma 33.47.3 says that the corresponding generic
section of $W_0$ is smooth. Its finite etale pullback is smooth as well,
and hence the geometrically irreducible pullback is geometrically
integral. The section has dimension $e-1$. The same argument applies
to the identity cover of $W_0$
\citep[Tags 0G4F and 0FD6, full proofs]{StacksDGenericSurface}.

Now choose a fixed tower $k=K_0\subset K_1\subset\cdots\subset K_{d-2}$
as follows: $K_i$ is an algebraic closure of
$K_{i-1}(t_{i,0},\ldots,t_{i,N})$, where each displayed set is
algebraically independent over $K_{i-1}$. Let $H_i$ be the hyperplane
with these coefficients. Successively intersect $W$ with these
hyperplanes and extend scalars at each step. The preceding paragraph,
applied to the identity cover, gives a smooth integral surface $W_*$
over $K=K_{d-2}$.

For every connected finite etale cover $V\to U$ defined over $k$, the
restriction $V_W\to W$ is connected: $V$ is normal and connected, hence
integral, and its inverse image of the dense open $W$ is nonempty and
integral. Apply the same paragraph at each successive stage to this
restriction. It follows that
\[
             V\times_U W_*\quad\hbox{is geometrically integral.}
             \tag{D.GS.2}
\]
Each cover may have its own nonempty Bertini open in a coefficient
space. Its generic point belongs to every such open. Thus the already
fixed tower works for all original covers simultaneously, without an
intersection of infinitely many open sets and without a countability
assumption. This is the geometric-generic quantifier also distinguished
in \citet[Remark 1.5 and Corollary 2.3, pp.~2 and 5--6]{Landesman2020D};
we have verified it directly for our successive hyperplanes.

Take the reduced closure $Z$ of $W_*$ in $Y_K$ and the reduced image
$C=q_K(Z)$ in $S_K$. The image is closed because $q_K$ is proper.
Both $Z$ and $C$ are integral surfaces, and $Z\to C$ is birational:
over the open $W_K$ the map $q_K$ is an isomorphism and the section
$W_*$ is closed in $W_K$. In particular,
\[
                        C\cap W_K=W_*.
\]
Let $X=C^\nu$. Normalization is finite for a scheme of finite type over
a field. It is an isomorphism over the normal open $W_*$. Hence $X$ is
a normal proper integral surface containing $W_*$ as a dense open,
and $g:X\to C\hookrightarrow S_K$ is finite
\citep[Tag 035E, Lemmas 29.55.8--11]{StacksDGenericSurface}.
Since $T_K$ is finite, its inverse image $T_X$ under this finite
morphism is finite as well. This is the step for which we use a finite
omitted set. Exceptional curves in $Z$ may be contracted in $C$; they
do not make $T_X$ positive-dimensional.

For a connected finite etale cover $V\to U$, its pullback to
$X\setminus T_X$ is finite etale over a normal integral base. Each of
its connected components is open and closed and maps surjectively to
the base. Every component therefore meets the dense open $W_*$.
Equation (D.GS.2) shows that this restriction is connected, so there is
only one component. The finite-cover criterion for surjectivity of
profinite fundamental groups proves (D.GS.1). The last assertion follows
because a quotient of a finitely generated profinite group is finitely
generated.
\end{proof}

\begin{corollary}[Proved: quantitative surface bounds at finite punctures in every dimension]
\label{D:normal-finite-puncture-quantitative-transfer}
Let $S$ be normal, proper and connected over an algebraically closed
field of characteristic $p>0$, let $\dim S\ge2$, and let $T\subset S$
be finite. Put $G=\pi_1(S\setminus T)$. There are constants $c,C$,
depending on this fixed pair, such that the following hold.

\begin{enumerate}
\item For every finite field $\Lambda$ and finite-dimensional continuous
$\Lambda[G]$-module $M$,
\[
                    \dim_\Lambda H^1(G,M)\le c\dim_\Lambda M.
\]
For each open subgroup $H\subset G$ the bound is
$c[G:H]\dim M$. The maximal prosolvable quotient of $G$ is finitely
generated.
\item For every $r,f\ge1$, with $q=p^f$,
\[
 \#\bigl(\operatorname{Hom}_{\rm cont}(G,\operatorname{GL}_r(\mathbf F_q))
                 /\operatorname{GL}_r(\mathbf F_q)\bigr)
                  \le q^{Cr^2}.
\]
\item Fix $B\ge1$, and put $q=p^f$ for $f\ge1$. Uniformly for finite subgroups
$\widetilde A\subset\operatorname{GL}_r(\mathbf F_q)$ whose scalar
kernel in the projective image $A$ has order at most $B$, the number
of homomorphisms $G\to A$, modulo the restricted ambient
$\operatorname{PGL}_r(\mathbf F_q)$-conjugacy relation, is at most
$q^{C_Br^2}$. In particular the simple-product conclusions proved for
these bounded-scalar families on surfaces hold for $G$ too. This
includes the stated defining-characteristic symplectic and orthogonal
families, and the linear and unitary families with bounded natural
centre.
\end{enumerate}
No degree-two cohomology bound is asserted by this transfer, and these
conclusions do not eliminate arbitrary nonabelian crowns.
\end{corollary}

\begin{proof}
Apply Theorem~\ref{D:generic-surface-finite-punctures} and write
$\Gamma=\pi_1(X\setminus T_X)\twoheadrightarrow G$ for its surjection.
The surface $X/K$ and its finite boundary are fixed, independently of
all coefficients, ranks, and finite covers. Inflation in degree one is
injective, so Theorem~\ref{D:punctured-surface-all-characteristics-H1}
gives (1). Shapiro's lemma gives its open-subgroup form. The maximal
prosolvable quotient is also a quotient of the finitely generated
maximal prosolvable quotient of $\Gamma$.

Precomposition with a surjection is injective on homomorphisms modulo
any fixed target-conjugacy relation: a conjugacy after precomposition
is already a conjugacy on $G$. Thus
Theorem~\ref{D:global-quadratic-representation-bound} gives (2), and
Proposition~\ref{D:bounded-scalar-projective-growth}, applied to the
punctured surface group as in
Corollaries~\ref{D:surface-bounded-centre-simple-product},
\ref{D:bounded-centre-unitary-simple-product} and
\ref{D:symplectic-orthogonal-simple-product}, gives (3).
Equivalently, every collection of distinct simple quotient kernels of
$G$ pulls back to such a collection for $\Gamma$, so its product group
is a quotient of the corresponding surface simple-product group.

There is no analogous unconditional injection on $H^2$ under this
surjection: the five-term sequence can have a nonzero preceding
contribution from $H^1(\ker(\Gamma\to G),M)^G$. The projective counts
were transferred directly, without such an injection.
\end{proof}
